\documentclass[border=10pt]{standalone}
\usepackage{fontspec}
\setmainfont{Noto Serif}[Ligatures=TeXOff]
\XeTeXgenerateactualtext=1
\newfontfamily\bengalifont{Noto Serif Bengali}[Script=Bengali, NFSSFamily=bnfam]
\usepackage[Latin,Bengali]{ucharclasses}
% group-free font switching (as in header-common.tex): robust inside braces, cells and quotes
\makeatletter
\def\BN@saved{\rmdefault}
\def\BN@bn{bnfam}
\newcommand\BNin{\edef\BN@tmp{\f@family}\ifx\BN@tmp\BN@bn\else\let\BN@saved\BN@tmp\fi\fontfamily{bnfam}\selectfont}
\newcommand\BNout{\fontfamily{\BN@saved}\selectfont}
\makeatother
\setTransitionsFor{Bengali}{\BNin}{\BNout}
% Typographic quotes and dashes are not in a ucharclasses class (class 0), so a Bangla word right
% after “ or ‘ would not trigger the Bangla font. Treat them as Latin so that every switch fires.
\XeTeXcharclass"2018=\XeTeXcharclass`A
\XeTeXcharclass"2019=\XeTeXcharclass`A
\XeTeXcharclass"201C=\XeTeXcharclass`A
\XeTeXcharclass"201D=\XeTeXcharclass`A
\XeTeXcharclass"2013=\XeTeXcharclass`A
\XeTeXcharclass"2014=\XeTeXcharclass`A
\XeTeXcharclass"2026=\XeTeXcharclass`A
\usepackage{newunicodechar}
\newunicodechar{।}{{\bengalifont\char"0964}}
\newunicodechar{॥}{{\bengalifont\char"0965}}

\usepackage[table]{xcolor}
\usepackage{booktabs,tabularx,array}
\definecolor{bapilrule}{HTML}{1F2937}
\definecolor{bapilstripe}{HTML}{F3F4F6}
\renewcommand{\arraystretch}{1.28}
\setlength{\tabcolsep}{5pt}

\begin{document}
\begin{minipage}{22cm}
{\large\bfseries সারণি ৫৯।}\ {\large ঝুঁকি-স্তর প্রশ্নমালা: স্কোরযুক্ত মাত্রা (সংযুক্তি খ)}\par\vspace{7pt}
\small
\rowcolors{2}{bapilstripe}{white}
\begin{tabularx}{\linewidth}{@{}>{\hsize=0.424\hsize\raggedright\arraybackslash}X>{\hsize=1.341\hsize\raggedright\arraybackslash}X>{\hsize=1.059\hsize\raggedright\arraybackslash}X>{\hsize=1.059\hsize\raggedright\arraybackslash}X>{\hsize=1.059\hsize\raggedright\arraybackslash}X>{\hsize=1.059\hsize\raggedright\arraybackslash}X@{}}
\arrayrulecolor{bapilrule}\toprule
\rowcolor{white}\textbf{কোড} & \textbf{মাত্রা} & \textbf{স্কোর ০} & \textbf{স্কোর ১} & \textbf{স্কোর ২} & \textbf{স্কোর ৩} \\
\midrule
D1 & সিদ্ধান্তে প্রভাব & শুধু গবেষণা & বিশ্লেষককে তথ্য দেয় & নীতি বা কর্মসূচি-সংক্রান্ত সরকারি সিদ্ধান্তে তথ্য দেয় & নির্দিষ্ট ব্যক্তি বা প্রতিষ্ঠান-সংক্রান্ত সিদ্ধান্তে তথ্য দেয় বা স্বয়ংক্রিয় করে \\
D2 & ডেটার সংবেদনশীলতা & P0 উন্মুক্ত & P1 অভ্যন্তরীণ, ব্যক্তিগত নয় & P2 ব্যক্তিগত & P3 সংবেদনশীল ব্যক্তিগত বা সংরক্ষিত \\
D3 & সম্ভাব্য প্রভাবিত মানুষ & কেউ নয় বা অভ্যন্তরীণ & ১০,০০০-এর কম & ১০,০০০ থেকে ১০ লক্ষ & ১০ লক্ষের বেশি \\
D4 & ক্ষতির প্রত্যাবর্তনযোগ্যতা & সহজেই প্রত্যাবর্তনযোগ্য & চেষ্টায় প্রত্যাবর্তনযোগ্য & প্রত্যাবর্তন কঠিন & অপ্রত্যাবর্তনযোগ্য \\
D5 & স্বয়ংক্রিয়তার মাত্রা & অনুরোধে পরামর্শ & প্রতিটি ফলাফল মানুষ পর্যালোচনা করেন & নমুনা মানুষ পর্যালোচনা করেন & সম্পূর্ণ স্বয়ংক্রিয় \\
D6 & প্রভাবিত গোষ্ঠীর ঝুঁকিপ্রবণতা & কোনো প্রভাবিত ব্যক্তি নেই & সাধারণ জনগণ & ঝুঁকিপূর্ণ গোষ্ঠী অন্তর্ভুক্ত & ঝুঁকিপূর্ণ গোষ্ঠীকে লক্ষ্য করে \\
D7 & উন্মুক্ততা & অভ্যন্তরীণ গবেষক & কর্মকর্তা & জনসাধারণের জন্য তথ্য & পরিণামসহ জনসাধারণের জন্য নির্দেশনা \\
\bottomrule
\end{tabularx}
\par\vspace{5pt}{\footnotesize\itshape স্তর-নিয়ম: কোনো নিষিদ্ধ ব্যবহার চিহ্ন থাকলে T0। অন্যথায় মোট স্কোর কমপক্ষে ১৩, বা D1 = ৩, বা D2 = ৩, বা D2 ও D3 উভয়ই কমপক্ষে ২ হলে T3; মোট কমপক্ষে ৬, বা D1 = ২, বা D7 কমপক্ষে ২ হলে T2; অন্যথায় T1। bapil.risk-এ বাস্তবায়িত।}
\par\vspace{3pt}{\scriptsize\color{gray} উৎস: BAPIL টেকনিক্যাল ব্লুপ্রিন্ট, সংস্করণ ১.০ (২০২৬)। CC BY 4.0।}
\end{minipage}
\end{document}
